\documentclass[12pt]{article}
\usepackage[utf8]{inputenc}
\usepackage[T1]{fontenc}
\usepackage{hyperref}
\title{Temiar Language Resources for K-12 Curriculum Development}
\author{Compiled by Claude Sonnet 5 (claude-sonnet-5) for Andr\'as Kornai}
\date{}
\begin{document}
\maketitle

\section{Introduction}
Temiar (ISO 639-3 tea; Glottocode temi1246) is an Austroasiatic language of the Orang Asli in Peninsular Malaysia. It had no standard written form for a long time. Linguists such as Benjamin and Diffloth adapted the IPA to record it, and it was largely unwritten apart from some printed pamphlets circulated by Baha'i religionists (per the encyclopedia entry).

\section{Primary Educational Resources}
No Temiar school primers or textbooks were found. Orang Asli education is discussed in a 2021 study of language education for Orang Asli children in Malaysia, which finds literacy well below the national average and argues for grounding language education in the children's own environment. The government programme Jom Belajar (2020) targets basic literacy and numeracy for indigenous children who have never attended school.
\begin{itemize}
\item \url{https://eprints.um.edu.my/35783}
\item \url{https://knova.um.edu.my/research_publications_2021_2025/10427/}
\item \url{https://www.bernama.com/en/news.php?id=2253293}
\end{itemize}

\section{Linguistic and Reference Materials}
\begin{itemize}
\item First Temiar dictionary, by seven Temiar collaborators with Nathalie Means: Temiar--English and English--Temiar, over 2,600 root words and over 3,800 derivatives: \url{https://temiar.com/dic.html}
\item Temiar.com main page: \url{https://temiar.com/temiar.html}
\item Orang Asli resources page: \url{https://temiar.com/oaresources.html}
\item Encyclopedia.com entry on the Temiar: \url{https://www.encyclopedia.com/humanities/encyclopedias-almanacs-transcripts-and-maps/temiar}
\end{itemize}

\section{Religious and Community Resources}
No Temiar Bible translation was found in these searches. A Temiar-related account appears at \url{https://borderlessjournal.com/tag/temiar} (a personal account of a Bible-study-based boarding school arrangement).

\section{Digital and Audio Resources}
None found beyond the temiar.com pages above.

\section{Curriculum Development Context}
The Temiar dictionary is the main published language resource found, and some copies were planned for distribution in Malaysia to support adult literacy in Temiar. No Temiar-medium school materials were found. The contents of the temiar.com pages other than the dictionary were not examined.
\end{document}